\chapter[Schöne Erfahrungen aus dem Gemeindeleben]{Schöne Erfahrungen aus dem Gemeindeleben, die andern Gemeinden zur Ermunterung dienen können}
\quellseite{3}

Gar manche unserer Gemeinden befinden sich nach der einen oder andern Seite hin in schwerer, bedrängter Lage. In einzelnen Gemeinden herrscht Unfriede, Streit und Parteiwesen, die Glieder halten nicht zusammen und es will nicht vorwärts, sondern eher rückwärts gehen. Andere Gemeinden haben einen schweren äußerlichen Stand; sie sind umgeben von Ungläubigen und Falschgläubigen, die ihnen gram und feind sind; wenig oder gar kein äußerliches Wachsthum ist zu verzeichnen; die alten Glieder sterben weg, neue schließen sich nicht an, und nur mit größter Mühe und Noth läßt sich die Gemeinde noch zusammenhalten. Wieder andere Gemeinden haben beständig mit Logen und gottlosen Vereinen zu kämpfen, die in die Gemeinde eindringen wollen oder schon eingedrungen sind, und ein ernstes, entschiedenes Auftreten droht die ganze Gemeinde zu sprengen. Noch andere Gemeinden seufzen unter einer schweren Schuldenlast und sind doch kaum im Stande, die laufenden Ausgaben für den Gemeindehaushalt aufzubringen, geschweige daß sie die Schulden verringern könnten. Solchen Gemeinden wird es zur Ermunterung und zum Troste gereichen, wenn sie hören, wie der HErr in Schwestergemeinden, die auch bedrückt waren, wunderbar geholfen, das treue Festhalten am Wort und das entschiedene Bekenntniß der Wahrheit herrlich gesegnet und die Herzen seiner Christen willig gemacht hat, für sein Reich und seine Gemeinde reichlich zu opfern. Zwei Vorkommnisse sind uns kürzlich bekannt geworden, die die Leser des „Lutheraner“ mit Freude und zum Segen hören werden. Das erste ist uns von einem Gliede der Gemeinde mitgetheilt worden, das zweite von dem Pastor der Gemeinde. Das Gemeindeglied schreibt folgendermaßen:

„Das Jahr 1903 war für unsere Gemeinde in — —“ (Name und Ort der Gemeinde thut nichts zur Sache) „ein ereignißvolles
\quellseite{4}
Jahr und wird den Gliedern unserer Gemeinde unvergeßlich bleiben. Ein Bericht von dem, was wir erfahren haben, mag auch andern von Nutzen sein, oder doch dazu dienen, daß sie sich mit uns freuen und mit uns Gottes Güte preisen.

„In den ersten Monaten des verflossenen Jahres starb eine Anzahl bewährter Glieder, die viele Jahre hindurch ihre Güter und Gaben treulich in den Dienst der Gemeinde gestellt hatten. Dieser Verlust wurde schmerzlich empfunden. Gar manchen drängten sich die Fragen auf: Wie werden wir jetzt die nöthigen Mittel aufbringen, um den Gemeindhaushalt fortzuführen? Wie können wir daran denken, unter diesen Umständen unsere Schulden völlig abzutragen? Kein Gedanke aber lag uns allen ferner als der, gerade jetzt unter diesen Umständen wider das Logenwesen zu Felde zu ziehen. Wir wußten aus Erfahrung, daß ein Kampf wider das Logenwesen unsere Gemeinde in den innersten Fugen erschüttern würde. Im Laufe des Jahres war nämlich eine Anzahl unserer Glieder beiderlei Geschlechts leider in das gottlose Logenwesen verstrickt worden.

„Doch was geschah? So sorgfältig wir uns auch hüteten, unter den obwaltenden Verhältnissen einen Logenkampf herbeizuführen, so wurden wir doch gerade jetzt in den Kampf hineingetrieben. Der Teufel und seine Werkzeuge wähnten, jetzt sei der passende Zeitpunkt gekommen, unserer Gemeinde den Garaus zu machen. Die äußerliche Veranlassung zum Kampf war ein Begräbnißfall. Eine Frau, die Glied unserer Gemeinde und zugleich auch Glied des ‘Women’s Relief Corps’ war, starb. Ihr nicht zur Gemeinde gehörender Mann stellte das Ansinnen an unsern Pastor, er solle gemeinschaftlich mit der Loge beim Begräbniß amtiren. Darauf konnte dieser selbstverständlich nicht eingehen. Er handelte hierin ganz im Sinne der Gemeinde, aber keineswegs im Sinne der Logen. Diese faßten seine Weigerung als eine große Beleidigung auf. Die Kriegserklärung folgte. Der Rumor begann. Namentlich fühlte sich die Frau eines Zeitungsherausgebers, die früher unsere Schule besucht hatte und bei uns confirmirt worden war, berufen, für die Logen in die Schranken zu treten. Monate lang erschienen von Zeit zu Zeit längere Artikel in den Zeitungen, in denen Pastor und Gemeinde in der gehässigsten Weise angegriffen wurden. Selbst die Rechte, die uns die Landesverfassung verbürgt, wurden beanstandet. Wir schwiegen, unsere Gegner aber wurden um so herausfordernder. So ging es drei Monate lang fort.

„Der Gräberschmückungstag rückte heran. Da erschien eines Tages eine Einladung des hiesigen deutschen Methodistenpredigers in den Zeitungen, gerichtet an die ‘Grand Army of the Republic’, die ‘Sons of Veterans’ und das ‘Women’s Relief Corps’, ihre Gedächtnisfeier in seiner Kirche abzuhalten. Dieser Einladung war die Bemerkung hinzugefügt, in seiner Kirche würden sie nicht beleidigt werden. Die Einladung wurde angenommen, der Methodistenpastor, ein abtrünniger Lutheraner, hielt die ‚Predigt‘ und ließ Tags darauf diese sogenannte Predigt durch den Druck veröffentlichen. Er mißbrauchte Gottes Wort, vergötterte die alten Soldaten und schmeichelte den Logen. Ihren Gipfelpunkt aber fanden seine Auslassungen darin, daß er wie mit Fingern auf unsern Pastor und unsere Gemeinde wies als auf Leute, die des Schutzes der Flagge nicht würdig und werth seien. Jetzt erkannte unsere Gemeinde, daß der Zeitpunkt gekommen sei, wo auch sie reden und ein ganz entschiedenes Zeugniß ablegen müsse. Dies geschah öffentlich durch den Druck sowohl von Seiten der Gemeinde als auch von Seiten des Pastors. Wir forderten Respectirung unserer Rechte, traten entschieden ein für unsern Pastor und bekannten unsere Stellung zu den Logen in Ausdrücken, die nicht mißverstanden werden konnten. Unser Zeugniß gipfelte in der Erklärung: Personen, die der gottfeindlichen Loge angehören, können nicht Glieder unserer Gemeinde werden, und diejenigen Glieder unserer Gemeinde, die noch den Logen angehören, müssen, nachdem so viel Belehrung über die falsche Religion und das widergöttliche Wesen der Logen stattgefunden hat, sich einmal entscheiden und entweder die Loge verlassen oder in Zucht genommen werden wegen ihres Ungehorsams gegen Gottes Wort.

„Diese unsere Erklärung erregte nicht geringes Aufsehen, und das allgemeine Urtheil war: die Gemeinde kann bei dieser Stellung zu den Logen nicht fortbestehen. Einige verstiegen sich so weit, daß sie meinten, die schöne Kirche der Gemeinde könne bald als Tanzsaal benutzt werden. Doch was geschah? Von höheren Logenbeamten erging die Aufforderung an die Herausgeber der Zeitungen, in diesem Handel nichts mehr zu veröffentlichen. So bekamen wir Ruhe von außen. Und welchen Ausgang nahm dieser Handel mit den Logengliedern innerhalb der Gemeinde? Wir durften die Freude erleben, daß mehrere von diesen erklärten: Wir sagen uns los von der Loge und bleiben bei der Gemeinde. Andere aber, mit denen wir schon vorher wegen Verachtung der Gnadenmittel hatten handeln müssen, ließen in der Logensache überhaupt nicht mehr mit sich verhandeln, sondern schieden aus. Mit andern verhandelt die Gemeinde noch jetzt. Kurz, die Scheidung hat stattgefunden.

„Wie steht es aber mit der Gemeinde selbst? Geht die über diesem Handel wirklich zu Grunde? Nimmermehr! Wohl scheiden solche aus, die schon längst innerlich nicht von uns gewesen sind. Bei der großen Menge der Gemeindeglieder aber hat der Kampf nur gute Wirkungen gehabt. Wir haben gelernt, aufs Wort zu merken. Wir sind schärfer geschieden von denen, die draußen sind, die Glieder der Gemeinde aber treten einander näher. Es ist, als ob ein Bann von uns genommen ist. Während die Gleichgültigkeit sich früher stark bemerkbar machte, fangen wir jetzt an, uns recht zu freuen, daß wir Lutheraner sind und als solche gewürdigt werden, die Wahrheit des Wortes Gottes zu bekennen. An Stelle der Ausscheidenden treten beständig neue Glieder ein, und zwar solche, die Liebe zum Wort befunden. Kurz, wir haben nichts verloren, wir haben nur gewonnen.

„Schließlich möchte noch gefragt werden: Wie steht es jetzt mit der finanziellen Lage der Gemeinde? Noch vor wenigen Jahren lastete eine Schuld von mehr als \$3000 auf uns. Der Rest dieser Schuld ist im verflossenen Jahre nicht nur gänzlich getilgt worden, sondern wir haben noch einen Ueberschuß von \$300 in der Kasse.

„‚Der Herr hat geholfen.‘ Dies erkennend, beschloß die Gemeinde, am Sonntag, den 10. Januar, einen besonderen Lob- und Dankgottesdienst abzuhalten. Wir haben gejubelt und in der Stille Gott gelobt zu Zion. Ja, ‚der Herr hat Großes an uns gethan, deß sind wir fröhlich‘.“

In einem Briefe bemerkt der Schreiber noch Folgendes: „Der Grund, warum ich diese Mittheilung (in völligen Einverständniß mit P. —) mache, ist der, daß dadurch andere möchten
\quellseite{5}
ermuntert werden, Stellung zu nehmen gegen das götzendienerische Logenwesen. Es würde aber auch unsern Brüdern hier in der Umgebung zur Stärkung dienen. Wo Kriege geführt und Schlachten geliefert werden, da gibt es Schwache und Verwundete, die der sorgfältigsten Pflege bedürfen. So steht es auch bei uns.“

Das zweite Beispiel, das andern Gemeinden zur Ermunterung dienen kann, ist dieses:

Vor Kurzem war im „Lutheraner“ von einer jungen Gemeinde berichtet worden, deren Glieder sich finanziell angestrengt und nicht nur ihre laufenden Ausgaben bestritten, sondern auch ihre Kirchbauschulden abgetragen hatten. Das war einem Leser unseres Blattes in einer andern Gemeinde aufgefallen, und er wandte sich an die Redaktion mit der Frage, ob denn auch alle diese Gelder in rechter Weise zusammengebracht worden seien, nämlich durch freiwillige Gaben, und nicht etwa durch Bazaare, fairs, oyster suppers, ice cream socials, Lotterien etc. Obwohl der Artikel selbst schon darauf hinwies und wir auch keinen Zweifel daran trugen, daß die Glieder in der rechten Weise, durch die Barmherzigkeit Gottes ermahnt und durch die Liebe Christi bewogen, ihre irdischen Gaben dargebracht hatten, wandten wir uns doch an den Pastor der Gemeinde, der uns Folgendes darüber schreibt:

„Das Geld ist auf rechte Weise zusammengekommen. Es ist unter uns noch nie Geld aufgebracht worden für kirchliche Zwecke außer durch freiwillige Beiträge. Das deutet ja auch der Bericht an, der vom ‚Collectiren‘ unserer Vorsteher, von einem ‚zweimaligen Collectiren unter unsern Gliedern‘, daß das Geld durch einmaliges Collectiren nicht zusammenkam, redet. Daher ist es kaum nöthig, daß der ‚Lutheraner‘ noch mehr darüber schreibt. Doch hinzufügen möchte ich noch, daß der ‚Lutheraner‘ nächst Gott ein Verdienst bei der Sache hat. Wir haben nämlich, wie seiner Zeit der ‚Lutheraner‘ das vorschlug und mit Gründen bewies, daß das monatliche Schulgeld abgeschafft und der Gehalt für den Lehrer mit auf die allgemeine Beitragsliste gesetzt, und das hat solche gute Frucht gebracht, daß wir im letzten Jahr von ungefähr denselben Leuten \$500 mehr einkommen haben für unsere Gemeindekasse und in Folge dessen im Stande waren, unsere Schule von 95 Kindern in zwei Klassen zu theilen und eine zweite Kraft zu berufen, während wir vorher kaum einen Lehrer erhalten konnten. Das ist gewiß ein Segen und auch nachahmenswerth für andere Gemeinden.“

Wer wollte sich nicht über solche schöne Erfahrungen in unsern Gemeinden freuen? Und ganz gewiß gibt es noch viele, viele andere Gemeinden unter uns, die auch in der einen oder andern Hinsicht auf besondere Weise die Hülfe und den Segen des Herrn erfahren haben. Kurze Worte, die manche Pastoren der Einsendung ihres Parochialberichtes für das „Statistische Jahrbuch“ beifügen, deuten darauf hin.

\begin{verse}
Der HErr ist noch und nimmer nicht\\
Von seinem Volk geschieden,\\
Er bleibet ihre Zuversicht,\\
Ihr Segen, Heil und Frieden.\\
Mit Mutterhänden leitet er\\
Die Seinen stetig hin und her.\\
Gebt unserm Gott die Ehre!
\end{verse}

L. F.
